\documentclass[11pt,letterpaper]{article}
\usepackage[lmargin=1in,rmargin=1in,bmargin=1in,tmargin=1in]{geometry}
\usepackage{style}

\pagenumbering{gobble}

% -------------------
% Content
% -------------------
\begin{document}

% TItle
\begin{center} 
\bfseries
\color{scred}
\LARGE Syllabus Quick Facts \par\vspace{0.2\baselineskip}
\Large MATH 142: Calculus II --- Fall 2026
\end{center} \pspace


% Course Information
\mysection{0.27}{Course Information}{course_info}
\hspace{0.53cm} {\itshape Instructor Email}: \href{mailto:cm264@mailbox.sc.edu}{cm264@mailbox.sc.edu} \par
\hspace{0.53cm} {\itshape Course Webpage}: \href{https://coffeeintotheorems.com}{https://coffeeintotheorems.com} \par
\hspace{0.53cm} {\itshape Office Hours}: The instructor's office is LeConte 345C. The office hours are XXXX \pspace


% Grading Components
\mysection{0.27}{Grading Components}{grade_comp}
Course grades are determined by the following components: \par \vspace{-0.3cm}
	\begin{table}[!ht]
        \begin{tabular}{lr}
	Check-Ins & 5\% \\
	Labs & 10\% \\
	Gateway Exams& 10\% \\
	Homework & 20\% \\
	Exam I -- III & 30\% \\
	Final Exam & 25\%
        \end{tabular} 
        \end{table}


% Attendance 
\mysection{0.27}{Attendance}{attendance}
Attend each lecture and show up on time. Address any absences---anticipated or otherwise---with the instructor. If you miss a lecture, you are responsible for any material covered, any work assigned, any course changes made, etc. during the class. Five or more unexcused absences from lectures could result in receiving a grade penalty per additional absence or an `F' in the course. Furthermore, excessive lateness will also count as an absence. 
\pspace


% Check-Ins 
\mysection{0.27}{Check-Ins}{check}
There will be a check-ins \textit{every} class, typically at the start of class. Because solutions will often then be immediately discussed, no make-up check-ins will be given (except under extraordinary circumstances). \pspace


% Labs 
\mysection{0.27}{Labs}{labs}
On various weeks, students will have a lab to complete in a Python Jupyter Notebook. These labs will help engage them with the material and learn some basic programming skills. Students will have a designated lab time to work on these labs. However, if a student does not complete their lab during this time, they are still expected to complete and submit the lab on time.
\pspace


% Gateways 
\mysection{0.27}{Gateways}{gateways}
There are two Gateway exams during the semester. The Gateway exams are 30~minute exams that will help students to achieve mastery over basic Precalculus/Calculus (Gateway~I) and integration skills (Gateway~II) and help assure students that they are prepared for the future material. 
\pspace



\newpage



% Homeworks 
\mysection{0.27}{Homeworks}{homeworks}
There will be weekly homework assignments. Homeworks will mostly be given and submitted using MyOpenMath. This assessment system is free for students and is integrated into Blackboard. Students will not need to create an account or make any purchases. 
\pspace


% Exams 
\mysection{0.27}{Exams}{exams}
There will be three exams in this course, each worth 10\% of the course grade, for a total of 30\% of the course grade. There will also be a final exam worth 25\% of the course grade. Together, all exams are worth 55\% of the course grade. 
\pspace


% Course Schedule 
\mysection{0.27}{Course Schedule}{schedule}
The following is a \emph{tentative} schedule for the course and is subject to change. 
        \begin{table}[!ht]
        \centering
        \scalebox{1}{%
        \begin{tabular}{ll || ll}
        Date & Topic(s) & Date & Topic(s) \\ \hline         
	08/18 & Integration Review & 10/13 & Absolute/Conditional Convergence \\
	08/19 & Integration Review & 10/14 & Recitation \\ 
	08/20 & Integration-by-Parts & 10/15 & Fall Break (No Classes) \\
	08/24 & Gateway I & 10/19 & Lab 5 \\
	08/25 & Integration-by-Parts & 10/20 & Ratio \& Root Test \\
	08/26 & Lab 1 & 10/21 & Recitation \\
	08/27 & Recitation & 10/22 & Power Series \\
	08/31 & Lab 2 & 10/26 & Lab 6 \\
	09/01 & Trigonometric Integrals & 10/27 & Exam 2 \\
	09/02 & Recitation & 10/28 & Recitation \\
	09/03 & Trig. Substitution & 10/29 & Taylor Series \\
	09/07 & Labor Day (No Classes) & 11/02 & Lab 7 \\
	09/08 & Partial Fractions & 11/03 & Election Day (No Classes) \\
	09/09 & Recitation & 11/04 & Lab 8 \& 9 \\
	09/10 & Partial Fractions & 11/05 & Taylor Series \\
	09/14 & Recitation & 11/09 & Recitation \\
	09/15 & Improper Integrals & 11/10 & Taylor Remainder Theorem \\
	09/16 & Recitation & 11/11 & Recitation \\
	09/17 & Applications \& Review & 11/12 & Series Applications \\
	09/21 & Recitation & 11/16 & Recitation \\
	09/22 & Exam 1 & 11/17 & Exam Review \\
	09/23 & Lab 3 & 11/18 & Recitation \\
	09/24 & Sequences \& Series & 11/19 & Exam 3 \\
	09/28 & Lab 4 & 11/23 & Thanksgiving Break (No Classes) \\
	09/29 & Geometric \& Alternating Series & 11/24 & Thanksgiving Break (No Classes) \\
	09/30 & Recitation & 11/25 & Thanksgiving Break (No Classes) \\
	10/01 & Telescoping Series \& Integral Test & 11/26 & Thanksgiving Break (No Classes) \\
	10/05 & Recitation & 11/30 & Recitation \\
	10/06 & Limit Comparison Test & 12/01 & Polar Coordinates \\
	10/07 & Recitation & 12/02 & Recitation \\
	10/08 & Direct Comparison Test & 12/03 & Parametric Functions \\
	10/12 & Recitation & 12/-- & Final Exam
        \end{tabular}
        }
        \end{table}

\end{document}